\documentclass[a4paper]{article}
\usepackage[utf8]{inputenc}
\usepackage[T1]{fontenc}
\usepackage[margin=1.5cm]{geometry}
\usepackage{tikz}
\usetikzlibrary{shapes.geometric}
\usepackage{adjustbox}

\tikzset{
  entita/.style={draw, rectangle, thick, minimum width=2.4cm, minimum height=0.9cm, fill=white},
  rel/.style={draw, diamond, aspect=2, thick, inner sep=1pt, fill=white},
  attr/.style={draw, circle, inner sep=0pt, minimum size=0.18cm, fill=white},
  idn/.style={draw, circle, inner sep=0pt, minimum size=0.18cm, fill=black},
  card/.style={font=\small, inner sep=2pt}
}

\begin{document}

\begin{center}
\begin{adjustbox}{max width=\textwidth}
\begin{tikzpicture}

  % ---------- ENTITA' ----------
  \node[entita] (Utente)      at (0,0)   {Utente};
  \node[entita] (Recensione)  at (7,0)   {Recensione};
  \node[entita] (Ricetta)     at (14,0)  {Ricetta};
  \node[entita] (TipoCibo)    at (14,6)  {TipoCibo};
  \node[entita] (Ingrediente) at (14,-6) {Ingrediente};

  % ---------- RELAZIONI ----------
  \node[rel] (Scrive)   at (3.5,0)  {Scrive};
  \node[rel] (Riguarda) at (10.5,0) {Riguarda};
  \node[rel] (HaTipo)   at (14,3)   {HaTipo};
  \node[rel] (Contiene) at (14,-3)  {Contiene};

  % ---------- COLLEGAMENTI + CARDINALITA' ----------
  \draw (Utente) -- node[card, above, pos=0.3] {(0,N)} (Scrive);
  \draw (Scrive) -- node[card, above, pos=0.7] {(1,1)} (Recensione);

  \draw (Recensione) -- node[card, above, pos=0.3] {(1,1)} (Riguarda);
  \draw (Riguarda)   -- node[card, above, pos=0.7] {(0,N)} (Ricetta);

  \draw (Ricetta) -- node[card, right, pos=0.3] {(1,N)} (HaTipo);
  \draw (HaTipo)  -- node[card, right, pos=0.7] {(0,N)} (TipoCibo);

  \draw (Ricetta)  -- node[card, right, pos=0.3] {(1,N)} (Contiene);
  \draw (Contiene) -- node[card, right, pos=0.7] {(0,N)} (Ingrediente);

  % ---------- ATTRIBUTI: Utente ----------
  \node[idn,  label={[font=\small]above:nomeUtente}] (u1) at (-1.6,1.6) {};
  \node[attr, label={[font=\small]above:email}]      (u2) at (0,1.9)    {};
  \node[attr, label={[font=\small]above:password}]   (u3) at (1.6,1.6)  {};
  \draw (Utente) -- (u1);  \draw (Utente) -- (u2);  \draw (Utente) -- (u3);

  % ---------- ATTRIBUTI: Recensione ----------
  \node[idn,  label={[font=\small]above:id}]    (r1) at (5.6,1.6) {};
  \node[attr, label={[font=\small]above:voto}]  (r2) at (7,1.9)   {};
  \node[attr, label={[font=\small]above:data}]  (r3) at (8.4,1.6) {};
  \node[attr, label={[font=\small]below:testo}] (r4) at (7,-1.7)  {};
  \draw (Recensione) -- (r1);  \draw (Recensione) -- (r2);
  \draw (Recensione) -- (r3);  \draw (Recensione) -- (r4);

  % ---------- ATTRIBUTI: Ricetta ----------
  \node[idn,  label={[font=\small]right:id}]                 (c1) at (16.4,1.6)  {};
  \node[attr, label={[font=\small]right:titolo}]             (c2) at (16.4,0.8)  {};
  \node[attr, label={[font=\small]right:istruzioni}]         (c3) at (16.4,0)    {};
  \node[attr, label={[font=\small]right:tempoPreparazione}]  (c4) at (16.4,-0.8) {};
  \node[attr, label={[font=\small]right:difficolta}]         (c5) at (16.4,-1.6) {};
  \foreach \i in {1,...,5} { \draw (Ricetta) -- (c\i); }

  % ---------- ATTRIBUTI: TipoCibo ----------
  \node[idn, label={[font=\small]right:nome}] (t1) at (16.2,6) {};
  \draw (TipoCibo) -- (t1);

  % ---------- ATTRIBUTI: Ingrediente ----------
  \node[idn, label={[font=\small]right:nome}] (i1) at (16.2,-6) {};
  \draw (Ingrediente) -- (i1);

  % ---------- ATTRIBUTI della relazione Contiene ----------
  \node[attr, label={[font=\small]left:quantita}] (q1) at (11.2,-2.4) {};
  \node[attr, label={[font=\small]left:unita}]    (q2) at (11.2,-3.6) {};
  \draw (Contiene) -- (q1);  \draw (Contiene) -- (q2);

\end{tikzpicture}
\end{adjustbox}
\end{center}

\end{document}